\documentclass[12pt]{article}

\input{../../_assets/preambulo.tex}


\begin{document}

    % 1. Foto de fondo
    % 2. Título
    % 3. Encabezado Izquierdo
    % 4. Color de fondo
    % 5. Coord x del titulo
    % 6. Coord y del titulo
    % 7. Fecha

    
    \input{../../_assets/portada}
    \portadaExamen{ffccA4.jpg}{Ecuaciones\\Diferenciales II\\Examen XIV}{Ecuaciones Diferenciales II. Examen XIV}{MidnightBlue}{-8}{28}{2026}{}

    \begin{description}
        \item[Asignatura] Ecuaciones Diferenciales II.
        \item[Curso Académico] 2022/2023.
        \item[Grado] Doble Grado en Ingeniería Informática y Matemáticas.
        \item[Grupo] Único.
        \item[Profesor] Rafael Ortega Ríos.
        \item[Descripción] Segunda prueba de clase.
        \item[Fecha] 31 de Mayo de 2023.
        % \item[Duración] ---.
    
    \end{description}
    \newpage


    % ------------------------------------
    
    \begin{ejercicio}
        ¿Existe una sucesión de funciones $\{f_n\}_{n \geq 1}$, $f_n : \left[0, 1\right] \to \bb{R}$, que sea equicontinua y uniformemente acotada pero que no converja puntualmente?
    \end{ejercicio}

    \begin{ejercicio}
        Se considera el problema de valores iniciales
        \[
        \ddot{x} + \operatorname{sen} x = 0, \quad x(0) = \veps, \quad \dot{x}(0) = 0
        \]
        y se denota por $x(t, \veps)$ a la solución. Calcula, si existe, $\deld{x}{\veps}(t, \pi)$.
    \end{ejercicio}

    \begin{ejercicio}
        La solución del sistema
        \[
        \dot{x} = y^2, \quad \dot{y} = x - 1, \quad x(0) = x_0, \quad y(0) = y_0
        \]
        se denota por $(x(t; x_0, y_0), y(t; x_0, y_0))$. Calcula $\deld{x}{y_0}(2; 1, 0)$.
    \end{ejercicio}

    \begin{ejercicio}
        Para cada $n \geq 1$ se denota por $x_n(t)$ a la solución del problema de valores iniciales
        \[
        \dot{x} = \frac{\operatorname{sen}(nx)}{n}, \quad x(0) = 2.
        \]
        Para cada $t \in \bb{R}$ calcula, si existe, $\lim\limits_{n \to \infty} x_n(t)$.
    \end{ejercicio}

    \begin{ejercicio}
        Se define la sucesión de funciones $\{f_n\}_{n \geq 0}$, $f_n : \left[0, 1\right] \to \bb{R}$, por el procedimiento iterativo
        \[
        f_0(t) = -1, \quad f_{n+1}(t) = 2 + \int_0^t \cos(f_n(s))ds, \quad t \in \left[0, 1\right].
        \]
        Demuestra que esta sucesión converge uniformemente en $\left[0, 1\right]$.
    \end{ejercicio}

    \newpage
    \setcounter{ejercicio}{0} % Reiniciar contador de ejercicios
    % \noindent
    % \textbf{Solución.}

\end{document}